% Part: first-order-logic
% Chapter: beyond
% Section: modal-logics

\documentclass[../../../include/open-logic-section]{subfiles}

\begin{document}

\olfileid{fol}{byd}{mod}

\olsection{مودالي منطقونه}

د يوې شرطي جملې دا بېلګه وګورئ:
\begin{quote}
  که جېرمي په هغې کوټه کښې يوازې وي، نو هغه مست او بربنډ دے او پر
  څوکيو ګډېږي۔
\end{quote}
دا داسې شرطي وينا ده چې ښايي په مادي لحاظ رښتونې خو بيا هم
غولوونکې وي، ځکه عبارت داسې ښيي چې د مقدم او پايلې تر منځ تر دې
پياوړې اړيکه شته چې يوازې يا مقدم دروغ وي يا تالي رښتونے۔ يعنې، عبارت
دا ښيي چې دعوه يوازې په دې ځانګړې نړۍ کښې رښتونې نۀ ده (چرته چې
ښايي ځکه په تش ډول رښتونې وي چې جېرمي په کوټه کښې يوازې نۀ دے)، بلکې
که مقدم رښتونے \emph{وے} نو پايله به هم رښتونې \emph{وه}۔ په نورو
ټکو، دا وينا داسې اخيستل کېدے شي چې دعوه يوازې په دې نړۍ کښې نۀ، بلکې
په هره ``ممکنه'' نړۍ کښې رښتونې ده؛ يا دا چې د دې ځانګړې نړۍ د محض
رښتياوالي پر خلاف، \emph{لازماً} رښتونې ده۔

مودالي منطق د دې ډول لزوم د معنا کولو دپاره جوړ شوے۔ مودالي بياني
منطق له عادي بياني منطق څخه د بکس عملګر په زياتولو ترلاسه کېږي؛ يعنې
که $!A$ !!a{formula} وي، نو $\Box !A$ هم دے۔ په وجداني ډول $\Box !A$
وايي چې $!A$ \emph{لازماً} رښتونے، يا په هره ممکنه نړۍ کښې رښتونے
دے۔ $\Diamond !A$ په عام ډول د $\lnot \Box \lnot !A$ لنډون ګڼل کېږي
او داسې لوستل کېږي چې $!A$ \emph{ممکن} رښتونے دے۔ البته موداليت په
پريديکاتي منطق کښې هم ورزياتېدلے شي۔

کرېپکي !!{structure}s مودالي منطق ته معنٰی پوهنه ورکولے شي؛ په حقيقت
کښې کرېپکي دا معنٰی پوهنه لومړے د مودالي منطق په پام کښې نيولو سره
جوړه کړې وه۔ دلته پر جزوي ترتيبونو د محدودېدو پر ځاے په عام ډول د
``ممکنو نړۍګانو'' يو سټ~$P$ او د نړۍګانو تر منځ دوه ځاييزه
``د لاسرسي'' اړيکه $\Atom{R}{x,y}$ لرو۔ په وجداني ډول
$\Atom{R}{p,q}$ وايي چې نړۍ~$q$ له~$p$ سره سازګاره ده؛ يعنې که موږ
په نړۍ~$p$ کښې ``يو''، بايد دا امکان ومنو چې نړۍ به د~$q$ په څېر وه۔

مودالي منطق کله د ``مفهومي'' منطق په نامه يادېږي، د ``مصداقي'' منطق
پر خلاف۔ د کلاسیک منطق په څېر د يوه مصداقي منطق ټاکل شوې معنٰی پوهنه
يوازې يوې، ``واقعي'' نړۍ ته اشاره کوي؛ خو د يوه ``مفهومي'' منطق معنٰی
پوهنه پر يوې پراخې هستپوهنې تکيه کوي۔ د لزوم د !!{structure} کولو تر
څنګ، موداليت د نورو ژبنيو جوړښتونو د !!{structure} کولو دپاره هم
کارېدلے شي، او $\Box$ او $\Diamond$ د اړوند کارونې له مخې بيا تفسيرېږي۔
د بېلګې په توګه:
\begin{enumerate}
\item د اثبات‌وړتيا په منطق کښې $\Box !A$ داسې لوستل کېږي چې ``$!A$
  د اثبات وړ دے'' او $\Diamond !A$ داسې چې ``$!A$ سازګار دے''۔
\item په معرفتي منطق کښې $\Box !A$ داسې لوستل کېدے شي چې ``زه په
  $!A$ پوهېږم'' يا ``زه پر $!A$ باور لرم''۔
\item په زماني منطق کښې $\Box !A$ داسې لوستل کېدے شي چې ``$!A$ تل
  رښتونے وي'' او $\Diamond !A$ داسې چې ``$!A$ کله کله رښتونے وي''۔
\end{enumerate}

غواړو منطق په هغو قاعدو او بديهي اصولو پياوړے کړو چې له موداليت سره
کار لري۔ د بېلګې په توګه، نظام~$\Log{S4}$ د بياني منطق عادي بديهي
اصول او قاعدې، او ورسره لاندې بديهي اصول لري:
\begin{align*}
& \Box (!A \lif !B) \lif (\Box !A \lif \Box !B)\\
& \Box !A \lif !A \\
& \Box !A \lif \Box \Box !A
\intertext{او دا قاعده هم: ``له $!A$ څخه $\Box !A$ پايله واخلئ''۔
  $\Log{S5}$ لاندې بديهي اصل ورزياتوي:}
& \Diamond !A \lif \Box \Diamond !A
\end{align*}
د دې بديهي اصولو بدل ډولونه ښايي د بېلابېلو کارونو دپاره مناسب وي؛ د
بېلګې په توګه، S5 په عام ډول د منطقي لزوم د مفهوم ځانګړوونکے ګڼل
کېږي۔ ښه خبره دا ده چې اکثره داسې معنٰی پوهنه موندلے شو چې د کرېپکي
په !!{structure}s کښې د لاسرسي اړيکه په طبيعي لارو محدودوي او د هغې
په نسبت د !!{derivation} نظام سم او بشپړ وي۔ د بېلګې په توګه،
$\Log{S4}$ د کرېپکي د هغو !!{structure}s له ټولګي سره سمون لري چې د
لاسرسي اړيکه ئې انعکاسي او انتقالي وي۔ $\Log{S5}$ د کرېپکي د هغو
!!{structure}s له ټولګي سره سمون لري چې د لاسرسي اړيکه ئې {\em عمومي}
وي، يعنې هره نړۍ له هرې بلې نړۍ څخه د لاسرسي وړ وي؛ نو $\Box !A$ هله
او يوازې هله صدق کوي چې $!A$ په هره نړۍ کښې صدق وکړي۔

\end{document}
